\documentclass[
%% TIKZ_CLASSOPTION %%
tikz
]{standalone}
\usepackage{amsmath}
\usetikzlibrary{matrix}
%% EXTRA_TIKZ_PREAMBLE_CODE %%
\begin{document}
%% TIKZ_CODE %%
\usetikzlibrary{positioning}
\definecolor{black}{HTML}{000000}
\tikzset{
    > = stealth,
    every node/.append style = {
        draw = none
    },
    every path/.append style = {
        arrows = ->,
        draw = black,
        fill = none
    },
    hidden/.style = {
        draw = black,
        shape = circle,
        inner sep = .75pt
    }
    manifest/.style = {
        draw = black,
        shape = square,
        inner sep = .75pt
    }
}
\tikz{
    \node[square] (x1) at (0,6) {$X_1$};
    \node[square] (x2) at (2,6) {$X_2$};
    \node[square] (x3) at (4,6) {$X_3$};
    
    \node[hidden] (wx1) at (0,4) {$W-X_1$};
    \node[hidden] (wx2) at (2,4) {$W-X_2$};
    \node[hidden] (wx3) at (4,4) {$W-X_3$};
    
    \node[hidden] (wy1) at (0,2) {$W-Y_1$};
    \node[hidden] (wy2) at (2,2) {$W-Y_2$};
    \node[hidden] (wy3) at (4,2) {$W-Y_3$};
    
    \node[square] (y1) at (0,0) {$Y_1$};
    \node[square] (y2) at (2,0) {$Y_2$};
    \node[square] (y3) at (4,0) {$Y_3$};
    
    \path (wx1) edge (wx2);
    \path (wx2) edge (wx3);
    
    \path (wx1) edge (wy2);
    \path (wx2) edge (wy3);
    
    \path (wy1) edge (wy2);
    \path (wy2) edge (wy3);
    
    \path (wy1) edge (wx2);
    \path (wy2) edge (wx3);
    
  }
\end{document}
